\section{Whittaker 范畴}
\markboth{Whittaker 范畴}{}

延续经典的类比，我们现在想构造一个范畴，它是
\[
  N(K)\backslash G(\A)/G(\mathbb O)
\]
上函数空间的几何对应物。这个范畴将记作 $D(\Bun\rt_N)$。此外，我们将定义一个满子范畴
\[
  \Whit(G)_{glob}\subset D(\Bun\rt_N),
\]
它是 $N(K)\backslash G(\A)/G(\mathbb O)$ 上函数空间中「对 $N(\A)$ 的一个固定非退化特征作等变」所得子空间的几何对应物。

这些范畴的构造用到 Ran 空间上的范畴这一形式，见 [GL:Ran]。

\ob{5.1}{blk:5.1}\textbf{空间。}考虑如下 ind-叠，记作 $\Bun_N^{polar}$，它生活在 $\mathrm{Ran}_X$ 之上：对有限集 $I$，令 $(\Bun_N^{polar})_{X^I}$ 分类五元组 $(\underline x,P_G,P_T,\kappa,\gamma)$，其中 $\underline x\in X^I$，$(P_G,P_T,\kappa)$ 是允许在 $\underline x$ 处有极点的 Plücker 数据，$\gamma$ 是 $T$-丛的同构
\[
  P_T\simeq\rho(\omega_X).
\]

对应
\[
  I\mapsto D\bigl((\Bun_N^{polar})_{X^I}\bigr)
\]
自然地是 Ran 空间上的一个范畴，我们记作 $D\bigl((\Bun_N^{polar})_{X^{fSet}}\bigr)$。它自然地带有单位结构与增广。我们把所得的范畴 $D\bigl((\Bun_N^{polar})_{\mathrm{Ran}_X,un}\bigr)$ 记作 $D(\Bun\rt_N)$。

\ob{5.2}{blk:5.2}\textbf{Whittaker 子范畴。}对每个 $I$，存在某个幂幺群胚 $(\mathcal H_N)_{X^I}$ 作用在 $(\Bun_N^{polar})_{X^I}$ 上（见 [Ga2]；事实上 $(\mathcal H_N)_{X^I}$ 是作用于 $\Bun_B$ 的 $\mathcal H_N$ 的一个版本），它带有一个特征
\[
  \chi : (\mathcal H_N)_{X^I}\to\Ga .
\]

我们把 $\Whit(G)_{X^I}$ 定义为相应的等变范畴
\[
  \Whit(G)_{X^I} := D\bigl((\Bun_N^{polar})_{X^I}\bigr)^{((\mathcal H_N)_{X^I},\chi)}.
\]

$(\mathcal H_N)_{X^I}$ 的幂幺性蕴含遗忘函子
\[
  \oblv_{\Whit,X^I} : \Whit(G)_{X^I}\to D\bigl((\Bun_N^{polar})_{X^I}\bigr)
\]
全忠实。该函子有右伴随，记作 $\ind_{\Whit,X^I}$。

\ob{5.2.1}{blk:5.2.1}我们把对应 $I\mapsto\Whit(G)_{X^I}$ 看作 Ran 空间上的一个范畴，简记作 $\Whit(G)$。它自然地是单位的、带增广的。我们将考虑相应的范畴
\[
  \Whit(G)_{\mathrm{Ran}_X}
  \quad\text{与}\quad
  \Whit(G)_{\mathrm{Ran}_X,un}.
\]
我们还将使用记号 $\Whit(G)_{glob}:=\Whit(G)_{\mathrm{Ran}_X,un}$，以及
\[
  \ind_{\Whit,glob} : D(\Bun\rt_N)\rightleftarrows \Whit(G)_{glob} : \oblv_{\Whit,glob}
\]
记相应的伴随函子对。

注意，范畴 $\Whit(G)_\emptyset$（定义为 $D(\Bun_N)$ 中关于相应群胚的等变对象范畴）典范等价于 $\mathrm{Vect}$。我们把 $\Whit(G)_\emptyset\subset\Whit(G)_{glob}$ 的相应单位对象记作 $\one_{\Whit(G)_{glob}}$。

\ob{5.3}{blk:5.3}\textbf{Casselman-Shalika 公式。}考虑由对称幺半范畴 $\Rep(\Ch G)$ 在 Ran 空间上搭起来的范畴；我们把它记作 $\Rep(\Ch G,X^{fSet})$。

\ob{5.3.1}{blk:5.3.1}相应的范畴 $\Rep(\Ch G,X^I)$ 可以如下显式描述。把 $\Ch G$ 看作 $X$ 上常值的群 D-概形，并令
\[
  \Ch G_{X^{fSet}} : I\mapsto \Ch G_{X^I}
\]
为 Ran 空间上相应的分解群 D-概形。

我们有：$\Rep(\Ch G,X^I)$ 是 $\Ch G_{X^I}$ 在 $X^I$ 上的 D-模（与联络相容者）上的表示范畴。

\ob{5.3.2}{blk:5.3.2}下述与 Casselman-Shalika 公式本质等价：

\begin{othm}{拟定理}{5.3.3}{thm:5.3.3}
存在 $\mathrm{Ran}_X$ 上范畴的自然等价
\[
  \Whit(G)\simeq\Rep(\Ch G,X^{fSet}).
\]
\end{othm}

\begin{annote}
「与 Casselman-Shalika 公式本质等价」是说：经典公式把 Whittaker 函数在某个群元素上的取值写成 $\check G$ 的表示的特征标（Schur 函数）在相应参数处的值；几何侧把「Whittaker 函数空间」换成范畴 $\Whit(G)$，把「$\check G$ 的表示」换成 $\Rep(\Ch G,X^{fSet})$，而「取值」这一操作则升级为范畴等价。\par\vspace{0.15em}{\raggedleft 【来源：原书 p.9 §5.3.2（A）；Casselman-Shalika 公式的经典表述（A）】\par}
\end{annote}

\ob{5.3.4}{blk:5.3.4}由上述定理，我们推出等价
\[
  \Whit(G)_{\mathrm{Ran}_X}\simeq\Rep(\Ch G,\mathrm{Ran}_X)
  \quad\text{与}\quad
  \Whit(G)_{\mathrm{Ran}_X,un}\simeq\Rep(\Ch G,\mathrm{Ran}_X,un).
\]
这里我们把 $\Rep(\Ch G,X^{fSet})$ 看作单位的、带增广的，依据是分解（群）D-概形 $I\mapsto\Ch G_{X^I}$ 上的余单位结构。

\ob{5.3.5}{blk:5.3.5}\emph{与 $\QCoh(\LocSys_{\Ch G})$ 的关系。}范畴 $\Rep(\Ch G,\mathrm{Ran}_X)$ 对我们的意义由下述事实说明。回忆有伴随函子对
\[
  \Loc_{\Ch G,spec} : \Rep(\Ch G,\mathrm{Ran}_X)\rightleftarrows \QCoh(\LocSys_{\Ch G}) : \Gamma_{\Ch G,spec}.
\]
我们有：

\begin{othm}{定理}{5.3.6}{thm:5.3.6}
函子 $\Gamma_{\Ch G,spec}$ 全忠实。
\end{othm}

\ob{5.3.7}{blk:5.3.7}容易看出，函子 $\Gamma_{\Ch G,spec}$ 的像落在满子范畴
\[
  \Rep(\Ch G,\mathrm{Ran}_X,un)\subset\Rep(\Ch G,\mathrm{Ran}_X)
\]
中。因此，函子 $\Loc_{\Ch G,spec}$ 经过局部化
\[
  \Rep(\Ch G,\mathrm{Ran}_X)\to\Rep(\Ch G,\mathrm{Ran}_X,un)
\]
分解。

\begin{annote}
定理 5.3.6 与 5.3.7 合起来给出 §7 反复使用的「谱侧落脚点」：$\Whit(G)_{glob}\simeq\Rep(\Ch G,\mathrm{Ran}_X,un)$，而后者经 $\Gamma_{\Ch G,spec}$ 全忠实地坐在 $\QCoh(\LocSys_{\Ch G})$ 里。于是关于 Whittaker 范畴的任何陈述，都可以翻译成谱侧关于拟凝聚层的陈述。\par\vspace{0.15em}{\raggedleft 【来源：原书 p.10 §5.3.5--5.3.7（A）】\par}
\end{annote}
